% §17 Exercises
%
% Cross-reference resolution status:
%
% DEFINED here:
%   \label{ex:ch17-1} through \label{ex:ch17-11}
%
% RESOLVED (used here, defined elsewhere):
%   \Cref{def:cipher-map}                     -- §17.1
%   \Cref{sec:bernoulli-error-model}          -- §17.2
%   \Cref{thm:cipher-composition}             -- §17.2
%   \Cref{sec:trapdoor-paradigm}              -- §17.3
%   \Cref{sec:fhe-sidebar}                    -- §17.4
%   \Cref{sec:boolean-trapdoor}               -- §17.5
%   \Cref{def:trapdoor-set}                   -- §17.5
%   \Cref{sec:maximizing-confidentiality}     -- §17.6
%   \Cref{def:ppv}                            -- chapter 7 PPV definition
%   \Cref{def:npv}                            -- chapter 7 NPV definition
%   \Cref{def:perfect-filter-line}            -- chapter 6 §6.4
%   \Cref{sec:cancellation}                   -- chapter 6 §6.3
%   \Cref{thm:composition}                    -- chapter 3 §3.3
%   \Cref{thm:kronecker}                      -- chapter 4 §4.4
%   \Cref{ax:element-indep}                   -- chapter 4 Axiom 1
%   \Cref{sec:bernoulli-data-type}            -- chapter 13 algebraic types
%
% 11 exercises, 3 star / 5 starstar / 3 starstarstar:
% - star x 3: cipher-map PPV/C numerical; composition of channel
%   matrices; FHE vs cipher map in two sentences
% - starstar x 5: multi-query leakage via Kronecker;
%   adaptive attack at O(|A|/(1-fprate)); confidentiality vs
%   space; conjunctive Boolean fprate; cancellation in trapdoor
%   complement-union
% - starstarstar x 3: PPV-based confidentiality bound derivation;
%   trapdoor set as Bernoulli<T> instance; Shannon entropy
%   confidentiality in bits per query

% Difficulty-marker macros (matching chapters 1-16):
\providecommand{\exstar}{$\star$}
\providecommand{\exstarstar}{$\star\star$}
\providecommand{\exstarstarstar}{$\star\star\star$}

\section*{Exercises}
\addcontentsline{toc}{section}{Exercises}

\begin{enumerate}

% -----------------------------------------------------------------------
% WARM-UP  (star x 3): exercises 1, 2, 3
% -----------------------------------------------------------------------

\item \exstar
\label{ex:ch17-1}
\textbf{Cipher-map PPV and confidentiality.}
A trapdoor set is configured with cipher-map rates $\fprate =
0.05$ and $\fnrate = 0$ on a workload with prevalence $\pi =
0.1$.
\begin{enumerate}
  \item[(a)] Substitute into \Cref{def:ppv} to compute $\PPV$.
    Verify your answer numerically.
  \item[(b)] Compute the confidentiality $C = 1 - \PPV$ from
    \Cref{sec:maximizing-confidentiality}. State in one sentence
    what $C$ means operationally: given a server $\top$-report,
    what fraction of the report's information is genuine and
    what fraction protects the plaintext?
\end{enumerate}

\item \exstar
\label{ex:ch17-2}
\textbf{Cipher-map composition.}
Let $Q_1$ and $Q_2$ be the $2 \times 2$ channel matrices of two
cipher maps from \Cref{sec:bernoulli-error-model}, both with
$\fnrate = 0$:
\[
  Q_i \;=\;
  \begin{pmatrix}
    1 - \fprate_i & \fprate_i \\
    0             & 1
  \end{pmatrix},
  \qquad i = 1, 2.
\]
Compute the composed matrix $Q_1 Q_2$ for the chained cipher map.
State the $(\fprate, \fnrate)$ point of the composed map and
verify that $\fnrate = 0$ is preserved under composition.

\item \exstar
\label{ex:ch17-3}
\textbf{FHE versus cipher map: when to use each.}
State in two sentences when a designer should choose FHE
(\Cref{sec:fhe-sidebar}) over a cipher-map approximation and
when the choice should go the other way. Name one application
that motivates each choice. (A two-sentence answer is
sufficient.)

% -----------------------------------------------------------------------
% MEDIUM  (starstar x 5): exercises 4, 5, 6, 7, 8
% -----------------------------------------------------------------------

\item \exstarstar
\label{ex:ch17-4}
\textbf{Multi-query leakage via Kronecker factorization.}
A trapdoor set answers a batch of $m = 10$ queries at
$\fprate = 0.05$ and $\fnrate = 0$, with the queries drawn
independently (chapter~4 Axiom~1 in force; see
\Cref{ax:element-indep}).
\begin{enumerate}
  \item[(a)] Apply the chapter~4 Kronecker theorem
    \Cref{thm:kronecker} to write the joint channel matrix for
    the batch as a Kronecker product of the per-query matrices.
    What is the shape of the joint matrix?
  \item[(b)] Bound the expected number of false-positive
    $\top$-reports across the batch. Compute the numerical
    expectation for $\fprate = 0.05, m = 10$.
  \item[(c)] State in one sentence why iid Kronecker
    factorization does not directly bound the joint
    confidentiality $C$: confidentiality is a posterior, and
    posteriors over multiple queries combine through Bayes'
    rule rather than through matrix products.
\end{enumerate}

\item \exstarstar
\label{ex:ch17-5}
\textbf{Adaptive recovery attack.}
Consider an adversary who queries a trapdoor set holding $|A| = n$
plaintext members at $\fprate = \epsilon$ and $\fnrate = 0$. The
adversary observes a $\top$-report on each query and adapts its
next query based on the report stream.
\begin{enumerate}
  \item[(a)] Construct an adaptive strategy that enumerates the
    plaintext domain in some prescribed order and, on each
    $\top$-report, records the queried element as a candidate
    member.
  \item[(b)] Argue that the expected number of queries needed to
    enumerate $A$ in full is at most $O(n / (1 - \epsilon))$.
    (Hint: each true positive is reported $1 - \fnrate = 1$
    times; each spurious $\top$-report contributes a false
    candidate the adversary discards via a second-pass
    verification.)
  \item[(c)] Discuss in two sentences how the adaptive setting
    relates to the per-query iid analysis: the per-query rates
    are unchanged, but the adversary's posterior over $A$
    grows monotonically with the query count.
\end{enumerate}

\item \exstarstar
\label{ex:ch17-6}
\textbf{Confidentiality versus space.}
Two cipher-map designs for the same plaintext set $A$: design~A
sets $\fprate = 0.01$ and design~B sets $\fprate = 0.001$, both
with $\fnrate = 0$. The workload has $\pi = 0.05$.
\begin{enumerate}
  \item[(a)] Compute $\PPV$ and the confidentiality
    $C = 1 - \PPV$ for both designs from
    \Cref{sec:maximizing-confidentiality}.
  \item[(b)] State which design is more confidential and quantify
    the gap.
  \item[(c)] Using the chapter~5 Bloom-filter space formula,
    state which design uses more storage per inserted element.
    Comment in one sentence on the tension: lower storage is
    the more confidential design.
\end{enumerate}

\item \exstarstar
\label{ex:ch17-7}
\textbf{Conjunctive Boolean leakage.}
A query $x \in A \cap B$ is answered by computing
$\tilde f_A(x) \wedge \tilde f_B(x)$ on two trapdoor sets
(\Cref{sec:boolean-trapdoor}), each at $\fprate = 0.05$ and
$\fnrate = 0$, with independent hash families.
\begin{enumerate}
  \item[(a)] Compute the overall $\fprate$ of the conjunction
    using the chapter~3 composition theorem
    \Cref{thm:composition}.
  \item[(b)] At workload prevalence $\pi = 0.01$ for the
    intersection $A \cap B$, compute the conjunction's $\PPV$
    and confidentiality.
  \item[(c)] Compare with the single-set case at the same $\pi$
    and the same per-set $\fprate$: which is more
    confidential? State in one sentence why.
\end{enumerate}

\item \exstarstar
\label{ex:ch17-8}
\textbf{Cancellation in a trapdoor-set complement-union.}
Apply the chapter~6 cancellation analysis of
\Cref{sec:cancellation} to a trapdoor-set expression of the
form $\complement(\tilde A \cup \tilde B)$, where both
$\tilde A$ and $\tilde B$ have $\fprate = 0.05$ and
$\fnrate = 0$.
\begin{enumerate}
  \item[(a)] By De Morgan, rewrite the expression as
    $\complement \tilde A \cap \complement \tilde B$. State
    the $(\fprate, \fnrate)$ point of each complemented
    operand. (Hint: complementing swaps the two error axes.)
  \item[(b)] Compute the $(\fprate, \fnrate)$ point of the
    intersection of the two complements using the chapter~3
    composition theorem.
  \item[(c)] State whether the resulting cipher map sits on
    the perfect-filter line of \Cref{def:perfect-filter-line}
    or escapes it. (Cancellation lets composed expressions
    leave the line; the question is which axis the resulting
    point sits on, if any.)
\end{enumerate}

% -----------------------------------------------------------------------
% CHALLENGING  (starstarstar x 3): exercises 9, 10, 11
% -----------------------------------------------------------------------

\item \exstarstarstar
\label{ex:ch17-9}
\textbf{Tight lower bound on confidentiality.}
Derive a tight lower bound on the confidentiality
$C(\fprate, \fnrate, \pi)$ of a cipher map, given the
parameters $\fprate, \fnrate, \pi \in [0, 1]$.
\begin{enumerate}
  \item[(a)] Write $C$ as a closed-form rational function of the
    three parameters, starting from the chapter~7 PPV formula
    \Cref{def:ppv} and the identity $C = 1 - \PPV$.
  \item[(b)] At fixed $\pi$, characterize the locus of
    $(\fprate, \fnrate)$ pairs that achieve a target
    confidentiality $C \ge C^{*}$ for some prescribed
    $C^{*} \in (0, 1)$. Show that the locus is a half-plane
    in the Bernoulli square cut by a line through the origin.
  \item[(c)] Specialize to $\fnrate = 0$ (perfect-filter line
    case) and give an explicit lower bound on $\fprate$ as a
    function of $\pi$ and $C^{*}$. Verify your bound on the
    \Cref{sec:maximizing-confidentiality} worked example
    ($\pi = 0.01, C^{*} = 0.832$ requires
    $\fprate \ge \sim 0.05$).
\end{enumerate}

\item \exstarstarstar
\label{ex:ch17-10}
\textbf{Trapdoor set as a $\Bernoulli{T}$ instance.}
Cast a trapdoor set (\Cref{def:trapdoor-set}) as an instance of
the algebraic-types framework of chapter~13.
\begin{enumerate}
  \item[(a)] Identify the underlying plaintext type $T$. For a
    Bloom-filter-backed trapdoor set on integers, what is
    $T$? Be explicit about whether the encryption layer
    changes $T$ or only its representation.
  \item[(b)] Write the $\Bernoulli{T}$ channel matrix that
    governs encrypted membership queries. State which
    Bernoulli-square point the trapdoor set occupies and how
    the encryption layer interacts (or fails to interact) with
    that point.
  \item[(c)] Compose with a second trapdoor set on the same
    type $T$ via the chapter~13 product construction and
    state the $(\fprate, \fnrate)$ point of the composed
    cipher map. Compare with the answer from
    \Cref{thm:cipher-composition}: the two routes should
    agree.
\end{enumerate}

\item \exstarstarstar
\label{ex:ch17-11}
\textbf{Information-theoretic confidentiality in bits per
query.}
Apply chapter~8's Shannon entropy machinery to compute the
information leakage of a cipher map on a Boolean query
``$x \in A$?'' in bits per query, as a function of
$\fprate, \fnrate, \pi$.
\begin{enumerate}
  \item[(a)] Model the adversary's per-query posterior as a
    Bernoulli random variable with parameter $\PPV$ for
    $\top$-reports and parameter $1 - \NPV$ for
    $\bot$-reports. Write the Shannon entropy of the prior
    (a Bernoulli at $\pi$) and the conditional entropies of
    the two posteriors.
  \item[(b)] The per-query leakage in bits is
    $L = H(\pi) - \mathbb{E}_{\text{report}}[H(\text{posterior})]$,
    the prior entropy minus the expected posterior entropy.
    Write $L$ as a closed-form function of
    $\fprate, \fnrate, \pi$.
  \item[(c)] Specialize to the \Cref{sec:maximizing-confidentiality}
    worked example ($\fprate = 0.05, \fnrate = 0, \pi = 0.01$)
    and compute $L$ numerically. Compare with the FHE limit
    ($\fprate = \fnrate = 0$, $L = H(\pi)$, the full
    prior-entropy leakage): the cipher map leaks a small
    fraction of $H(\pi)$ per query, and the framework gives
    a quantitative measure of that fraction.
\end{enumerate}

\end{enumerate}
